\section*{Capítulo \ref{ch:g11:seismic-structure} --- Las ondas sísmicas revelan el interior}

\begin{solution}{exo:g11:seismic-structure:1}
Una \omterm{def:g11:seismic-structure:body-waves}{onda P} comprime y estira la roca, que se mueve en la dirección de
propagación; una \omterm{def:g11:seismic-structure:body-waves}{onda S} la cizalla, y la roca se mueve perpendicularmente a
esa dirección. Las \omterm{def:g11:seismic-structure:body-waves}{ondas P} atraviesan sólidos y líquidos; las \omterm{def:g11:seismic-structure:body-waves}{ondas S}, solo
sólidos. La \omterm{def:g11:seismic-structure:body-waves}{onda P} se registra primero, porque es la más rápida de las dos.
\end{solution}

\begin{solution}{exo:g11:seismic-structure:2}
$40 \times 111.2 = \qty{4448}{km}$, unos \qty{4450}{km}. $8900/111.2
\approx \ang{80}$.
\end{solution}

\begin{solution}{exo:g11:seismic-structure:3}
\qty{456}{s} son 7~minutos 36~segundos; \qty{822}{s} son 13~minutos
42~segundos. El retraso S--P es de $822 - 456 = \qty{366}{s}$, es decir,
6~minutos 6~segundos.
\end{solution}

\begin{solution}{exo:g11:seismic-structure:4}
$\sin i_c = 6.0/8.0 = 0.75$, luego $i_c \approx \ang{48.6}$.
\end{solution}

\begin{solution}{exo:g11:seismic-structure:5}
El núcleo externo tiene $5150 - 2891 = \qty{2259}{km}$ de espesor. El radio
del núcleo es $6371 - 2891 = \qty{3480}{km}$, y $3480/6371 \approx 0.55$: algo
más de la mitad del radio de la Tierra.
\end{solution}

\begin{solution}{exo:g11:seismic-structure:6}
6~minutos 6~segundos son \qty{366}{s}: el terremoto está a \ang{40}, es decir,
a $40 \times 111.2 \approx \qty{4450}{km}$. Un solo retraso solo da la
distancia: el epicentro puede estar en cualquier punto de un círculo de ese
radio alrededor de la estación. Hacen falta otras dos estaciones para
encontrar el único punto que está a la distancia correcta de las tres.
\end{solution}

\begin{solution}{exo:g11:seismic-structure:7}
$i_c \approx \ang{48.6}$ y $\cos i_c \approx 0.661$. Onda directa: $200/6.0
\approx \qty{33.3}{s}$. \omterm{ex:g11:seismic-structure:sp}{Onda de cabeza}: $200/8.0 + 2 \times 35 \times 0.661/6.0
\approx 25.0 + 7.7 = \qty{32.7}{s}$. La \omterm{ex:g11:seismic-structure:sp}{onda de cabeza} llega primero, unos
\qty{0.6}{s} antes.
\end{solution}

\begin{solution}{exo:g11:seismic-structure:8}
$\tan i_c \approx 1.134$: la distancia crítica es $2 \times 35 \times 1.134
\approx \qty{79}{km}$. La \omterm{prop:g11:seismic-structure:crossover}{distancia de cruce} es $2 \times 35 \times
\sqrt{14.0/2.0} = 70\sqrt 7 \approx \qty{185}{km}$. La estación del
\cref{exo:g11:seismic-structure:7} está a \qty{200}{km}, más allá del cruce:
la \omterm{ex:g11:seismic-structure:sp}{onda de cabeza} llega primero, como se encontró.
\end{solution}

\begin{solution}{exo:g11:seismic-structure:9}
$\sin i_c = 6.5/8.0 = 0.8125$, $i_c \approx \ang{54.3}$ y $\cos i_c \approx
0.583$. \omterm{prop:g11:seismic-structure:head-wave}{Tiempo de intercepción}: $2 \times 6 \times 0.583/6.5 \approx \qty{1.1}{s}$.
\omterm{prop:g11:seismic-structure:crossover}{Distancia de cruce}: $12\sqrt{14.5/1.5} \approx \qty{37}{km}$. Bajo un
continente, el cruce está hacia \qty{165}{km}, así que unas estaciones
separadas unas decenas de kilómetros captan las dos rectas; bajo el océano,
la \omterm{ex:g11:seismic-structure:sp}{onda de cabeza} ya llega primero más allá de unos \qty{37}{km}, y la recta
de la onda directa solo mide unas decenas de kilómetros: las estaciones
tienen que estar a pocos kilómetros unas de otras para trazarla.
\end{solution}

\begin{solution}{exo:g11:seismic-structure:10}
$\sin r = (8.06/13.7)\sin\ang{30} \approx 0.294$, luego $r \approx \ang{17.1}$:
el rayo se acerca a la perpendicular, es decir, hacia el centro de la Tierra.
Por eso atraviesa el núcleo por un camino largo cerca de su centro, y vuelve
a salir muy lejos, al otro lado: los rayos P que entran en el núcleo salen
mucho más lejos del terremoto que los rayos que lo rozan sin entrar.
\end{solution}

\begin{solution}{exo:g11:seismic-structure:11}
$12\,742/13.7 \approx \qty{930}{s}$, unos 15.5~minutos, mientras que la onda
tarda \qty{1210}{s}, unos 20~minutos: llega con unos 280~segundos de retraso.
Incluso a la velocidad más alta del manto la onda llegaría mucho antes, así
que el centro de la Tierra tiene que frenarla. Oldham concluyó que la Tierra
tiene un núcleo central en el que las \omterm{def:g11:seismic-structure:body-waves}{ondas P} viajan más despacio que en el
manto.
\end{solution}

\begin{solution}{exo:g11:seismic-structure:12}
(1) Por la proposición, $t_i = 2h\cos i_c/v_1$, y $\cos i_c = \sqrt{1 -
v_1^2/v_2^2} = \sqrt{v_2^2 - v_1^2}/v_2$. Así que
\[
  t_i = \frac{2h\sqrt{v_2^2 - v_1^2}}{v_1 v_2},
  \qquad\text{de donde}\qquad
  h = \frac{t_i\,v_1 v_2}{2\sqrt{v_2^2 - v_1^2}}.
\]
(2) $h = 6.23 \times 6.2 \times 8.1/\bigl(2\sqrt{65.61 -
38.44}\bigr) \approx 312.9/10.43 \approx \qty{30.0}{km}$.
\end{solution}

\begin{solution}{exo:g11:seismic-structure:13}
(1) $(3480/6371)^3 \approx 0.163$: el núcleo representa alrededor del
\qty{16}{\percent} del volumen. (2) La densidad media es la media ponderada
por los volúmenes: $5.52 = 0.163
\rho_c + 0.837 \times 4.45$, luego $\rho_c \approx (5.52 - 3.72)/0.163 \approx
\qty{11.0}{g/cm^3}$. (3) Es más que \qty{9.9}{g/cm^3}, como debe ser, porque
la densidad del núcleo aumenta hacia el centro (hasta \qty{13.1}{g/cm^3});
es unas cuatro veces la densidad de la corteza, \qty{2.6}{g/cm^3}. El núcleo
está hecho de un material mucho más denso que las rocas de la superficie.
\end{solution}

\begin{solution}{exo:g11:seismic-structure:14}
(1) $105 \times 111.2 \approx \qty{11700}{km}$ y $143 \times 111.2 \approx
\qty{15900}{km}$. (2) $(\cos\ang{105} - \cos\ang{143})/2 \approx (-0.259 +
0.799)/2 \approx 0.27$: alrededor del \qty{27}{\percent} de la superficie, un
anillo alrededor del punto opuesto al terremoto. (3) La \omterm{def:g11:seismic-structure:shadow-zone}{zona de sombra} es un
intervalo de distancias al epicentro: se desplaza con el terremoto. Cada
terremoto tiene su propio anillo.
\end{solution}

\begin{solution}{exo:g11:seismic-structure:15}
(1) $h = 9.5 \times 6.3 \times 8.1/\bigl(2\sqrt{65.61 - 39.69}\bigr) \approx
484.8/10.18 \approx \qty{48}{km}$. (2) $\sin i_c = 6.3/8.1 \approx 0.778$,
$\cos i_c \approx 0.629$, y $t_i = 2 \times 30 \times 0.629/6.3 \approx
\qty{6.0}{s}$. (3) La corteza bajo la cordillera tiene unos \qty{48}{km} de
espesor, más que la media continental de unos \qty{41}{km} y cerca de la
media de las cordilleras: las cordilleras altas se apoyan en una corteza
gruesa.
\end{solution}

\begin{solution}{pb:g11:seismic-structure:1}
\textbf{1.} Los puntos Pg están en una recta; los puntos Pn están en otra
recta, de menor pendiente.

\textbf{2.} El terremoto está a muy poca profundidad bajo el epicentro: una
estación a distancia cero registraría la onda al instante.

\textbf{3.} $(32.3 - 3.3)/(200 - 20) = 29.0/180 \approx \num{0.161}$~\unit{s}
por kilómetro.

\textbf{4.} $v_1 = 180/29.0 \approx \qty{6.2}{km/s}$.

\textbf{5.} $3.3 + 280 \times 0.161 \approx \qty{48.4}{s}$.

\textbf{6.} $(43.2 - 18.6)/(300 - 100) = 24.6/200 = \num{0.123}$~\unit{s}
por kilómetro, luego $v_2 = 200/24.6 \approx \qty{8.13}{km/s}$.

\textbf{7.} Está cerca de \qty{8.0}{km/s}: bajo la corteza está la roca de la
parte más alta del manto.

\textbf{8.} $\sin i_c = 6.21/8.13 \approx 0.763$, luego $i_c \approx
\ang{49.8}$.

\textbf{9.} La \omterm{ex:g11:seismic-structure:sp}{onda de cabeza} baja con el ángulo $i_c$ y vuelve a subir con
el mismo ángulo. El primer punto en el que puede volver a la superficie está
a $h\tan i_c$ más allá del punto en el que llegó al \omterm{def:g11:seismic-structure:moho}{Moho}, que a su vez está a
$h\tan i_c$ de la fuente: ninguna estación situada a menos de $2h\tan i_c$
puede recibirla.

\textbf{10.} Un sismómetro registra todas las ondas que le llegan, no solo la
primera; la onda Pn llega después de la onda Pg, pero se sigue viendo en el
registro como una onda separada.

\textbf{11.} A \qty{160}{km}, la Pg (\qty{25.8}{s}) llega antes que la Pn
(\qty{26.0}{s}); a \qty{180}{km}, la Pn (\qty{28.5}{s}) llega antes que la Pg
(\qty{29.0}{s}). El cambio se produce entre estas dos estaciones.

\textbf{12.} Pg: $t = 3.3 + 0.1611\,(x - 20) \approx 0.08 + 0.1611\,x$. Pn:
$t = 18.6 + 0.123\,(x - 100) = 6.30 + 0.123\,x$ ($x$ en kilómetros, $t$ en
segundos).

\textbf{13.} $0.08 + 0.1611\,x = 6.30 + 0.123\,x$ da $x \approx
6.22/0.0381 \approx \qty{163}{km}$.

\textbf{14.} La Pn, al cabo de $6.30 + 0.123 \times 250 \approx \qty{37.1}{s}$.

\textbf{15.} $t_i \approx \qty{6.3}{s}$.

\textbf{16.} $\cos i_c \approx 0.646$, luego $h = 6.21 \times 6.3/(2 \times
0.646) \approx \qty{30.3}{km}$.

\textbf{17.} $x_c = 2 \times 30.3 \times \sqrt{14.34/1.92} \approx 60.6 \times
2.73 \approx \qty{165}{km}$, a menos de \qty{2}{km} de los \qty{163}{km} de
la pregunta~13.

\textbf{18.} $\tan i_c \approx 1.18$: $2 \times 30.3 \times 1.18 \approx
\qty{72}{km}$. La estación más cercana que registra la Pn está a
\qty{80}{km}, y la de \qty{60}{km} no la registra: las dos cosas concuerdan.

\textbf{19.} Los tiempos se leen con \qty{0.1}{s} de precisión, y un error de
\qty{0.1}{s} solo en $t_i$ cambia $h$ en unos \qty{0.5}{km}; las velocidades,
sacadas de pendientes, también son inciertas, y el terremoto no estaba
exactamente en la superficie. Además, la corteza no es perfectamente
uniforme. Una precisión de uno o dos kilómetros es lo máximo que permiten
los datos.

\textbf{20.} \textbf{El \omterm{def:g11:seismic-structure:moho}{Moho} está a unos \qty{30}{km} bajo esta región}: la
corteza es aquí más bien delgada para un continente, cuya corteza tiene de
media unos \qty{41}{km}; \qty{30}{km} es lo típico de las regiones donde la
corteza se ha estirado.
\end{solution}
